\documentclass[11pt]{article}
\usepackage{fontspec,amsmath,amssymb,amsthm,hyperref}
\usepackage[margin=1in]{geometry}
\setlength{\emergencystretch}{3em}
\title{A reference for every initial coupling with one output law}
\author{LangevinLean mathematical exposition}
\date{}
\begin{document}\maketitle
\noindent Source module: \texttt{UnadjustedLangevin/ProductReference.lean}. Compile this proof document with XeLaTeX or LuaLaTeX. The accompanying HTML pairs the mathematical argument with the exact Lean statements and records the scope of the formalization.


This construction turns the analytic reference estimate into an exact initial Wasserstein term by realizing each finite initial coupling on a product with one fixed noise source.


\section{pointwiseUpdate}


Define the ordinary Euler update \(F_h(x,z)=x-h\nabla U(x)+\sqrt{2h}z\) in an arbitrary Hilbert-space presentation. It is the pointwise function whose lifted composition gives the \(L^2\) numerical chain.



\noindent Lean declarations: \texttt{pointwiseUpdate}.

\section{pointwiseUpdate measurable}


The Lipschitz gradient is continuous. Compose it with the first coordinate projection and combine the remaining continuous scalar and vector operations. The uncurried map \((x,z)\mapsto F_h(x,z)\) is therefore measurable.



\noindent Lean declarations: \texttt{pointwiseUpdate\_measurable}.

\section{productNoise memLp}


The noise projection from \(c\otimes Q\) to \(Q\) is measure preserving because the coupling is a probability. Pull back the supplied \(L^2(Q)\) noise function along this projection. Its square integrability is unchanged.



\noindent Lean declarations: \texttt{productNoise\_memLp}.

\section{productNoise}


Use the preceding square-integrability proof to form the class of \(((x,y),s)\mapsto W_n(s)\) on the concrete extension. Both the numerical chain and reference use this same \(L^2\) innovation.



\noindent Lean declarations: \texttt{productNoise}.

\section{productNoise coe}


The representative selected by \texttt{toLp} agrees almost everywhere with the pulled-back original noise \(W_n(s)\). This identity connects quotient-space noise arithmetic to the actual source random functions.



\noindent Lean declarations: \texttt{productNoise\_coe}.

\section{productOutputLaw}


Push \(\mu\otimes Q\) through the pointwise noise-driven chain starting at \(x\) with source path \(W(s)\). This definition omits the reference coupling entirely, making explicit the common algorithm output law for all coupling realizations.



\noindent Lean declarations: \texttt{productOutputLaw}.

\section{product chain coe}


Induct through the lifted chain and use the representative identities for initial state, product noise, lifted drift, addition, and scalar multiplication. At each fixed index the result agrees almost everywhere with the ordinary recursion driven by \(x\) and \(W(s)\).



\noindent Lean declarations: \texttt{product\_chain\_coe}.

\section{product chain outputLaw}


Replace the \(L^2\) chain by its almost-everywhere pointwise recursion. That recursion depends only on the algorithm's initial coordinate and noise sample, so \texttt{initialExtension\_map} gives its coupling-independent law. The result is exactly \texttt{productOutputLaw}.



\noindent Lean declarations: \texttt{product\_chain\_outputLaw}.

\section{productProbability}


A coupling of the initial probability with the reference measure is itself a probability, and its product with the fixed probability noise source is a probability. This local instance supplies normalization on every extension separately.



\noindent Lean declarations: \texttt{productProbability}.

\section{ProductReference}


For every finite-cost initial coupling, the structure supplies a concrete \texttt{L2Reference} on that coupling's product extension, with its initial state matching the second coordinate. Algorithm output laws and initial costs are proved outside the structure; no convergence estimate is a field.



\noindent Lean declarations: \texttt{ProductReference}.

\section{ProductReference.output finite}


Finite initial Wasserstein distance supplies at least one finite-cost initial coupling. On its concrete extension, the numerical and reference variables belong to \(L^2\), giving a finite-cost output coupling. Their established laws identify its marginals with the common output law and \(\pi\), hence the output Wasserstein distance is finite.



\noindent Lean declarations: \texttt{ProductReference.output\_finite}.

\section{dalalyan small step of product reference}


For each finite-cost initial coupling \(c\), apply the analytic \(L^2\) Dalalyan theorem on \(c\otimes Q\). The initial class-distance theorem identifies its initial error exactly with the real coupling cost, and the common output-law theorem identifies its numerical law.



The resulting bound holds for every finite \(c\), with coefficient \((1-mh)^n\ge0\). Infimizing their costs gives the genuine \(W_2(\mu,\pi)\), without selecting an optimal coupling. Output finiteness is proved first so the final real-valued inequality cannot conceal infinity.



\noindent Lean declarations: \texttt{dalalyan\_small\_step\_of\_product\_reference}.

\section{dalalyan large step of product reference}


Repeat the couplingwise analytic argument for the stable large-step branch, whose contraction factor is \(Lh-1\in[0,1)\). The geometric error sum has denominator \(2-Lh>0\). The finite-coupling infimum theorem replaces the initial class distance by the actual initial Wasserstein distance and yields the displayed large-step bias term.



\noindent Lean declarations: \texttt{dalalyan\_large\_step\_of\_product\_reference}.

\end{document}
